\section{Materiały dodatkowe}

\subsection{Repozytorium projektu}
Kod źródłowy projektu jest dostępny w repozytorium:
\begin{center}
  \url{https://github.com/PiotrKosiara/Projekt_ADAC}
\end{center}

Repozytorium zawiera:
\begin{itemize}
  \item kod frontendu testbedowego,
  \item backend FastAPI wraz z modelem danych i endpointami,
  \item pipeline feature engineering i treningu modelu,
  \item bot runner Playwright,
  \item konfigurację Docker Compose,
  \item raport LaTeX oraz materiały prezentacyjne.
\end{itemize}

\subsection{Struktura datasetu}
Dane projektowe są przechowywane przede wszystkim w PostgreSQL. Logicznie dataset
można podzielić na cztery warstwy:
\begin{enumerate}
  \item \textbf{Warstwa sesji}: tabela \texttt{sessions}, czyli identyfikator sesji,
  etykieta referencyjna, źródło danych i fingerprint środowiska.
  \item \textbf{Warstwa surowych eventów}: tabela \texttt{raw\_events}, czyli
  chronologiczny zapis \texttt{mousemove}, \texttt{click}, \texttt{scroll},
  \texttt{focus}, \texttt{blur} i zdarzeń viewportu.
  \item \textbf{Warstwa cech}: okna zachowania generowane dynamicznie przez backend,
  z cechami takimi jak \texttt{acc\_max}, \texttt{speed\_std},
  \texttt{path\_length}, \texttt{event\_type\_entropy}.
  \item \textbf{Warstwa wyników}: tabele \texttt{predictions},
  \texttt{model\_versions} i \texttt{enforcement\_actions}.
\end{enumerate}

\subsection{Rozmiar i skład danych końcowych}
Końcowy zestaw danych wykorzystany w raporcie:
\begin{itemize}
  \item 202 oznaczone sesje,
  \item 101 sesji \texttt{human},
  \item 101 sesji \texttt{bot},
  \item 644 próbki okien zachowania w ewaluacji aktywnego modelu,
  \item scenariusze botowe: \texttt{linear} -- 25 sesji,
  \texttt{human\_like} -- 26 sesji, \texttt{adaptive} -- 50 sesji.
\end{itemize}

\subsection{Odtwarzalność eksperymentu}
Podstawowy eksperyment można odtworzyć lokalnie przez:
\begin{lstlisting}[style=reportcode]
docker compose up -d --build
docker compose exec bot_runner python main.py --scenario linear --sessions 25 --record-video
docker compose exec bot_runner python main.py --scenario human_like --sessions 26 --record-video
docker compose exec bot_runner python main.py --scenario adaptive --sessions 50 --record-video
docker compose exec backend python -m app.ml.train --window-size 120 --stride 80 --tree-max-depth 12 --tree-min-samples-leaf 2
\end{lstlisting}

Dane \texttt{human} wymagają ręcznych interakcji z frontendem i etykietowania
sesji jako \texttt{human}. Ten element jest celowo manualny, ponieważ stanowi
referencyjne zachowanie człowieka dla scenariuszy botowych.
